\clearpage
\VerdierPhysicalPageAnchor{0128}%
\begingroup
\small
\setlength{\emergencystretch}{1.5em}%
\setlength{\abovedisplayskip}{3pt plus 2pt minus 2pt}%
\setlength{\belowdisplayskip}{3pt plus 2pt minus 2pt}%

\VerdierSemanticAnchor{ch2-2}
\section{Localisation dans les catégories triangulées.}

\VerdierSemanticAnchor{ch2-2-1}
\subsection{Systèmes multiplicatifs de morphismes [16].}

\VerdierSemanticAnchor{ch2-2-1-1}
\noindent\textbf{2.1.1.}\quad
\VerdierSemanticAnchor{ch2-2-1-1-multiplicative-system}
Soit $C$ une catégorie (non nécessairement additive). Un ensemble $S$
de morphismes de $C$ est appelé un \emph{système multiplicatif} si il
possède les propriétés suivantes :

\VerdierSemanticAnchor{ch2-2-1-1-sm1}
\smallskip
\noindent\textbf{SM1)} Le composé de deux morphismes composables de $S$
est un élément de $S$. Pour tout objet $X$ de $C$, le morphisme identique
de $X$ est un élément de $S$.

\VerdierSemanticAnchor{ch2-2-1-1-sm2}
\smallskip
\noindent\textbf{SM2)} Tout diagramme :

\VerdierSemanticAnchor{ch2-2-1-1-sm2-right-input}
\begin{center}
  \begin{tikzcd}[column sep=4.1em,row sep=2.5em]
    & {} \arrow[d,"s\in S"] \\
    {} \arrow[r,"f"'] & {}
  \end{tikzcd}
\end{center}
\noindent peut se compléter en un diagramme commutatif :

\VerdierSemanticAnchor{ch2-2-1-1-sm2-right-completion}
\begin{center}
  \begin{tikzcd}[column sep=4.1em,row sep=2.5em]
    {} \arrow[r,dashed,"g"] \arrow[d,dashed,"t\in S"']
      & {} \arrow[d,"s\in S"] \\
    {} \arrow[r,"f"'] & {}
  \end{tikzcd}
  \quad .
\end{center}
\noindent Tout diagramme :

\VerdierSemanticAnchor{ch2-2-1-1-sm2-left-input}
\begin{center}
  \begin{tikzcd}[column sep=4.1em,row sep=2.5em]
    {} \arrow[r,"g"] \arrow[d,"t\in S"'] & {} \\
    {} & {}
  \end{tikzcd}
\end{center}
\noindent peut se compléter en un diagramme commutatif :

\VerdierSemanticAnchor{ch2-2-1-1-sm2-left-completion}
\begin{center}
  \begin{tikzcd}[column sep=4.1em,row sep=2.5em]
    {} \arrow[r,"g"] \arrow[d,"t\in S"']
      & {} \arrow[d,dashed,"s\in S"] \\
    {} \arrow[r,dashed,"f"'] & {}
  \end{tikzcd}
  \quad .
\end{center}

\VerdierSemanticAnchor{ch2-2-1-1-sm3-boundary}
\noindent\textbf{SM3)} Pour tout couple de morphismes
$f,g\colon X\rightrightarrows Y$, les propriétés suivantes sont
équivalentes :
\endgroup
